\chapter*{Zusammenfassung}
\addcontentsline{toc}{chapter}{Zusammenfassung}

Karotisplaque gilt als Marker der systemischen Atherosklerose, ist in Populationskohorten aber
nicht als Variable erfasst: Er liegt nur implizit vor, in Ultraschallbildern, die niemand
befundet hat. Jede Analyse, die Plaque mit molekularen Daten in Beziehung setzt, muss den
Phänotyp daher zunächst mit einem Modell erzeugen und ist durch die Güte dieses Schrittes
begrenzt. Die vorliegende Arbeit sucht nach Plasmaproteinen, die mit subklinischer
Karotisatherosklerose assoziiert sind, und untersucht die Messung, von der eine solche Suche
abhängt.

Ein publiziertes Detektionsmodell wurde auf den gesamten Karotis-Ultraschallbestand der UK
Biobank angewandt und lieferte Detektionen für 179{.}611 Bilder von 20{.}014 Teilnehmenden bei
einer Plaqueprävalenz von 43,0\,\% gegenüber den ursprünglich berichteten 45\,\%. Gegen das
Olink-Proteompanel geprüft, zeigte dieser Phänotyp in den 2{.}542 Teilnehmenden mit beiden
Messungen keine Assoziation, die die Korrektur für multiples Testen überstand: 151 von 2{.}923
Proteinen erreichten nominelle Signifikanz, wo 146 zufällig zu erwarten sind, und der beste
Klassifikator erzielte eine kreuzvalidierte ROC-AUC von 0,541 gegenüber 0,551 für Alter und
Geschlecht allein. Das Alter war in denselben Daten signifikant assoziiert, was belegt, dass das
Verfahren Assoziationen findet, wo welche vorliegen.

Da ein fehlklassifizierter Prädiktor Assoziationen proportional zum Youden-Index $J$ abschwächt,
wurde das Modell anschließend außerhalb seiner Entwicklungskohorte geprüft, auf einem
öffentlichen argentinischen Datensatz. Seine Spezifität fiel von 89,2\,\% auf 62,5\,\%, während
die Sensitivität lediglich von 89,5\,\% auf 83,3\,\% zurückging. Auf diesem Datensatz
feinjustierte Instanzsegmentierungsmodelle hoben $J$ von 0,458 auf 0,625, das Zusammenführen
zweier Modelle auf 0,667, wobei alle Arbeitspunkte auf den Validierungs- und nicht auf den
Testdaten gewählt wurden. Eine auf denselben Bildern trainierte Detektorkontrolle erreichte
dasselbe Niveau, womit die Verbesserung der Domänenanpassung und nicht der
Segmentierungsformulierung zuzuschreiben ist.

Wählt man die am besten passende vorhergesagte Instanz unabhängig von ihrem Konfidenzwert, steigt
der mittlere Dice-Koeffizient von 0,655 auf 0,794, und kein Bild bleibt ohne korrekte Maske: Die
Modelle umreißen Plaques, die sie nicht melden. Ursache ist die Gewichtung der Zielfunktion, in
der der Konfidenzterm fünfundzwanzigmal niedriger gewichtet war als die geometrischen Terme. Die
begrenzende Komponente ist somit die Entscheidung auf Bildebene, nicht die Segmentierungsgüte.
